\documentclass[11pt]{amsart}
\usepackage{amsmath,amssymb,amsthm}
\usepackage[margin=1in]{geometry}

\newtheorem{theorem}{Theorem}[section]
\newtheorem{proposition}[theorem]{Proposition}
\newtheorem{lemma}[theorem]{Lemma}
\newtheorem{corollary}[theorem]{Corollary}
\theoremstyle{remark}
\newtheorem{remark}[theorem]{Remark}

\newcommand{\ve}{\varepsilon}
\newcommand{\R}{\mathbb{R}}
\newcommand{\Z}{\mathbb{Z}}
\newcommand{\Lop}{\mathcal{L}}

\title{$\ve$-uniform convergence of the monotone scheme:\\ the travelling wave case (working notes)}
\date{\today}

\begin{document}
\maketitle

\noindent\textbf{Status.} Working notes. Sections~\ref{sec:setting}--\ref{sec:main}
contain a complete argument for the non-sharp estimate of
Theorem~\ref{thm:main}, modulo the profile estimates of Lemma~\ref{lem:profile},
for which we give the proof in outline. Section~\ref{sec:sharp} describes how to
reach the sharp rates suggested by the experiments.

% ---------------------------------------------------------------------------
\section{Setting}\label{sec:setting}

\subsection*{Equation} One space dimension, $w^\ve\equiv\ve w_0$ with $w_0>-1$
constant, so $Y=u+\ve w_0$ and
\begin{equation}\label{eq:P}
u_t=u_{xx}-g_\ve(u),\qquad g_\ve(s)=(s+\ve w_0)f_\ve(s),\quad f_\ve(s)=\ve^{-2}f(s/\ve).
\end{equation}
\textbf{(F)} $f$ is Lipschitz (constant $L_f$), nonincreasing, $f>0$ on
$(-\infty,1)$ and $f=0$ on $[1,\infty)$. \textbf{(F$'$)} In addition
$f\in C^{1,1}(\R)$.

\subsection*{Scheme} Uniform grid $x_j=jh$, $j\in\Z$, time step $\Delta t$,
$t_n=n\Delta t$, $\delta^2V_j=(V_{j+1}-2V_j+V_{j-1})/h^2$. For bounded
sequences $V=(V^n_j)$ define
\begin{equation}\label{eq:L}
\Lop[V]^n_j=\frac{V^{n+1}_j-V^n_j}{\Delta t}-\delta^2V^{n+1}_j+f_\ve(V^n_j)\,\big(V^{n+1}_j+\ve w_0\big).
\end{equation}
The scheme (S) is $\Lop[U]=0$ with $U^0$ given. (It is the scheme of the
paper with $W\equiv\ve w_0$; constants are discrete-caloric.)

\subsection*{Comparison on $\Z$} A bounded $V$ is a \emph{supersolution} if
$\Lop[V]\ge0$ and $V+\ve w_0\ge0$; a \emph{subsolution} if $\Lop[V]\le0$.

\begin{lemma}\label{lem:comp}
Assume (F). If $V$ is a bounded supersolution, $U$ the bounded solution of (S)
with $U^0+\ve w_0\ge0$, and $U^0\le V^0$, then $U^n\le V^n$ for all $n$. The
analogous statement holds for subsolutions.
\end{lemma}

\begin{proof}
As in the paper, $D=U^{n+1}-V^{n+1}$ satisfies
$(\Delta t^{-1}+F_j)D_j-\delta^2D_j\le0$ for all $j$, where
$F=f_\ve(U^n)\ge0$. Let $s=\sup_jD_j$ and assume $s>0$. Given $\eta>0$ pick $j$
with $D_j>s-\eta$. Then $-\delta^2D_j\ge\frac{2}{h^2}(D_j-s)\ge-\frac{2\eta}{h^2}$
and $(\Delta t^{-1}+F_j)D_j\ge\Delta t^{-1}(s-\eta)$, so
$0\ge\Delta t^{-1}(s-\eta)-2\eta h^{-2}$, which fails for small $\eta$.
\end{proof}

% ---------------------------------------------------------------------------
\section{Travelling waves of \eqref{eq:P}}

For $A>\ve$ let $\Phi=\Phi_\ve$ be the travelling wave: $u=\Phi(x-ct)$ with
$c=c_\ve$,
\begin{equation}\label{eq:tw}
\Phi''+c\Phi'=g_\ve(\Phi),\qquad \Phi(-\infty)=-\ve w_0,\quad\Phi(+\infty)=A,
\end{equation}
normalized by $\Phi(0)=\ve$. Then $\Phi(z)=A-(A-\ve)e^{-cz}$ for $z\ge0$, and
$\Phi(z)=\ve\varphi(z/\ve)$ for $z\le0$, where $\varphi$ solves
$\varphi''+\sigma\varphi'=(\varphi+w_0)f(\varphi)$, $\sigma=\ve c$ (see the
paper). Let $c_0=\sqrt{2M}/A$ and $\kappa_0=\sqrt{f(-w_0)}$.

\begin{lemma}[Profile estimates]\label{lem:profile}
Assume (F). There are $\ve_0>0$ and constants $K,\theta_0,\mu>0$ depending only
on $f,w_0,A$ such that for $0<\ve\le\ve_0$:
\begin{enumerate}
\item $|c_\ve-c_0|\le K\ve$ and $c_\ve\le 2c_0$;
\item $\Phi'>0$ on $\R$, $\|\Phi'\|_\infty\le K$, $\|\Phi''\|_\infty\le K/\ve$,
$\|\Phi'''\|_\infty\le K/\ve^2$; under (F$'$) also $\|\Phi''''\|_\infty\le K/\ve^3$;
on $[\ve,\infty)$ all derivatives up to order four are bounded by $K$;
\item for every $Z\ge1$ there is $\theta(Z)\ge\theta_0e^{-\kappa_0 Z}$ with
$\Phi'\ge\theta(Z)$ on $[-Z\ve,2\ve]$;
\item for $z\le0$: $0\le\Phi(z)+\ve w_0\le K\ve e^{\kappa_0 z/(2\ve)}$, and
$|\Phi^{(k)}(z)|\le K\ve^{1-k}e^{\kappa_0 z/(2\ve)}$, $k=1,2,3$;
\item $\sup_z|\Phi_\ve(z)-u_0(z)|\le K\ve$, where $u_0(z)=A(1-e^{-c_0z})^+$.
\end{enumerate}
\end{lemma}

\begin{proof}[Outline]
In the inner variables, $\varphi_\sigma$ is the orbit of the planar system
$\varphi'=p$, $p'=-\sigma p+(\varphi+w_0)f(\varphi)$ leaving the saddle
$(-w_0,0)$ (eigenvalues $\frac12(-\sigma\pm\sqrt{\sigma^2+4f(-w_0)})$) along its
unstable manifold and reaching $\varphi=1$ with $p=\sigma(A/\ve-1)$. As
$\sigma\to0$ it converges, with its derivatives, to the heteroclinic of the
conservative system $p^2/2=G(\varphi)$, uniformly on half lines $(-\infty,0]$;
this gives (2)--(4) on $z\le0$ (the decay rate is at least
$\kappa_0/2$ for $\sigma$ small). Item (1) is Proposition~2.1 of the paper
together with $\sigma=\ve c_\ve$. On $z\ge0$, $\Phi$ is explicit. For (5),
compare both explicit expressions on $z\ge0$ and use $-\ve w_0\le\Phi\le\ve$ on
$z\le0$. Monotonicity: $\Phi'>0$ on $z\le0$ because $p>0$ along the unstable
manifold until $\varphi=1$; on $z\ge0$ it is explicit. Item (3) follows from
the exponential behaviour near the saddle and positivity on compacts.
\end{proof}

% ---------------------------------------------------------------------------
\section{Barriers}

For $\xi,q\colon\{t_n\}\to[0,\infty)$ nondecreasing let
\[
V^n_j=\Phi\big(x_j-ct_n+\xi_n\big)+q_n,\qquad
\sigma_n=\frac{\xi_{n+1}-\xi_n}{\Delta t},\qquad \dot q_n=\frac{q_{n+1}-q_n}{\Delta t}.
\]
$V+\ve w_0\ge0$ since $\Phi\ge-\ve w_0$. Write $z'=x_j-ct_{n+1}+\xi_{n+1}$ and
$z=x_j-ct_n+\xi_n=z'+(c-\sigma_n)\Delta t$. We assume throughout
\begin{equation}\label{eq:small}
h\le\ve,\qquad c\Delta t\le\ve,\qquad \sigma_n\le c.
\end{equation}

\begin{lemma}[Residual]\label{lem:residual}
Assume (F) and \eqref{eq:small}. Then
\[
\Lop[V]^n_j\ \ge\ \dot q_n+\sigma_n\Phi'(z')-r_1-r_2+f_\ve(V^n_j)q_{n+1}-B,
\]
where, with $I=[z'-h,\,z'+\max(h,c\Delta t)]$,
\[
r_1=\tfrac12c^2\Delta t\,\sup_I|\Phi''|,\qquad
r_2=\begin{cases}\tfrac h3\sup_I|\Phi'''|&\text{under (F)},\\[2pt]
\tfrac{h^2}{12}\sup_I|\Phi''''|&\text{under (F$'$)},\end{cases}
\]
and $B=0$ if $\Phi(z')\ge\ve$, while in general
\[
0\le B\le\frac{L_f}{\ve^3}\big(\Phi(z')+\ve w_0\big)\big(c\Delta t\sup_I\Phi'+q_n\big).
\]
\end{lemma}

\begin{proof}
\emph{Time difference.} By Taylor,
$\Phi(z')-\Phi(z)=-(c-\sigma_n)\Delta t\,\Phi'(z')+R$ with
$|R|\le\frac12(c\Delta t)^2\sup_I|\Phi''|$, so
$(V^{n+1}_j-V^n_j)/\Delta t=-c\Phi'(z')+\sigma_n\Phi'(z')+\dot q_n+R/\Delta t$.

\emph{Space difference.}
$\delta^2\Phi(z')-\Phi''(z')=h^{-2}\int_{-h}^h(h-|s|)\big(\Phi''(z'+s)-\Phi''(z')\big)ds$,
which is bounded by $\frac h3\sup|\Phi'''|$; under (F$'$), symmetrizing and
expanding to the next order gives $\frac{h^2}{12}\sup|\Phi''''|$.

\emph{Reaction.} Since $z\ge z'$ and $\Phi$ is increasing,
$\Phi(z)+q_n\ge\Phi(z')$, so $f_\ve(V^n_j)\le f_\ve(\Phi(z'))$ and
\[
f_\ve(V^n_j)(V^{n+1}_j+\ve w_0)=g_\ve(\Phi(z'))+f_\ve(V^n_j)q_{n+1}
-\underbrace{\big(f_\ve(\Phi(z'))-f_\ve(V^n_j)\big)(\Phi(z')+\ve w_0)}_{=:B\ge0}.
\]
If $\Phi(z')\ge\ve$ both values of $f_\ve$ vanish. Otherwise
$|f_\ve(a)-f_\ve(b)|\le L_f\ve^{-3}|a-b|$ and
$0\le\Phi(z)+q_n-\Phi(z')\le c\Delta t\sup_I\Phi'+q_n$.

\emph{Collecting.} Using \eqref{eq:tw} at $z'$,
$-c\Phi'-\Phi''+g_\ve(\Phi)=0$.
\end{proof}

% ---------------------------------------------------------------------------
\section{Main estimate}\label{sec:main}

Let
\[
\rho=\frac{\Delta t}{\ve^2}+\frac{h}{\ve^2}\ \text{ under (F)},\qquad
\rho=\frac{\Delta t}{\ve^2}+\frac{h^2}{\ve^3}\ \text{ under (F$'$)}.
\]

\begin{theorem}\label{thm:main}
Assume (F) [resp. (F$'$)], $w_0>-1$ constant, $A>0$, and let $U$ be the solution
of (S) with $U^0_j=\Phi_\ve(x_j)$. For every $T>0$ there are $C,\gamma>0$,
depending only on $f,w_0,A,T$, such that if $\ve\le\ve_0$, $h\le\ve$,
$c_0\Delta t\le\ve/2$ and $\rho\le\gamma$, then
\[
\max_{t_n\le T}\ \sup_j\big|U^n_j-\Phi_\ve(x_j-c_\ve t_n)\big|\le C\rho,
\qquad
\max_{t_n\le T}\ \sup_j\big|U^n_j-u(x_j,t_n)\big|\le C(\ve+\rho),
\]
where $u(x,t)=A\big(1-e^{-c_0(x-c_0t)}\big)^+$ is the travelling wave of the
limit free boundary problem.
\end{theorem}

In particular, with $\Delta t\le\ve^3$ and $h\le\ve^3$ (Lipschitz $f$) or
$h\le\ve^2$ (under (F$'$)), the scheme converges to the limit problem with
rate $O(\ve)$, uniformly in $\ve$.

\begin{proof}
We prove the upper bound; the lower bound is symmetric (use
$\Phi(x-ct-\xi)-q$; in Lemma~\ref{lem:comp} for subsolutions one uses
$U+\ve w_0=Y\ge0$). $K,\kappa_0,\theta(\cdot)$ are as in
Lemma~\ref{lem:profile}; $C_1,C_2,\dots$ depend only on $f,w_0,A,T$. Let
$a=1$ under (F) and $a=2$ under (F$'$); note that
\begin{equation}\label{eq:rho}
\Delta t\le\ve^2\rho,\qquad h^a\le\ve^2\rho .
\end{equation}

\emph{Constants $\mu,\delta,Z$.} Let $\mu=f(-w_0)/4>0$ and choose $\delta>0$ with
$f(s)\ge2\mu$ for $s\le-w_0+3\delta$. Choose $Z\ge1$ with
$Ke^{\kappa_0/2}e^{-\kappa_0Z/2}\le\delta$ and $L_fKe^{-\kappa_0Z/2}\le\mu/4$; this fixes
$\theta=\theta(Z)>0$. With $\zeta=z'/\ve$ we split into the regions
(I) $\zeta\ge2$, (II) $-Z\le\zeta\le2$, (III) $\zeta\le-Z$.
Note $I\subset[z'-\ve,z'+\ve]$ by \eqref{eq:small}.

\emph{Residual bounds from Lemma~\ref{lem:profile}.}
\begin{itemize}
\item[(I)] $I\subset[\ve,\infty)$, where derivatives are $O(1)$:
$r_1+r_2\le C_1(\Delta t+h^a)$. Since $\Phi(z)+q_n\ge\Phi(z')\ge\ve$, $B=0$.
\item[(II)] $r_1\le C_1\Delta t/\ve\le C_1\ve\rho$ and $r_2\le C_1\rho$
(from $|\Phi'''|\le K\ve^{-2}$, resp. $|\Phi''''|\le K\ve^{-3}$). Since
$\Phi(z')+\ve w_0\le(1+w_0)\ve$, $B\le C_1\ve^{-2}(\Delta t+q_n)$.
\item[(III)] $r_1+r_2\le C_1\rho$ and, by item (4) and the choice of $Z$,
$B\le\frac{\mu}{4}\ve^{-2}(c\Delta tK+q_n)$.
\end{itemize}

\emph{Stability in (III).} There $z\le z'+\ve\le(1-Z)\ve\le0$, so by item (4)
$(\Phi(z)+\ve w_0)/\ve\le Ke^{\kappa_0(1-Z)/2}\le\delta$. If
$q_n\le\ve\delta$, then $V^n_j/\ve\le-w_0+2\delta$, so $f_\ve(V^n_j)\ge2\mu\ve^{-2}$
and $f_\ve(V^n_j)q_{n+1}\ge2\mu\ve^{-2}q_n$.

\emph{Choice of $q$ and $\xi$.} Set
\[
q_n=q_*+C_2(\Delta t+h^a)t_n,\qquad q_*=C_3\ve^2\rho,\qquad
\xi_n=\sigma t_n,\quad\sigma=C_4\rho .
\]
By \eqref{eq:rho}, $q_n\le(C_3+C_2T)\ve^2\rho$ for $t_n\le T$. If $\gamma$ is
small, $\rho\le\gamma$ implies $q_n\le\ve\delta$ and $\sigma\le c_0/2\le c$, so
\eqref{eq:small} and the stability bound hold.

\emph{Sign of $\Lop[V]$} (Lemma~\ref{lem:residual}).
\begin{itemize}
\item[(I)] $\Lop[V]\ge\dot q_n-r_1-r_2\ge(C_2-C_1)(\Delta t+h^a)\ge0$ if $C_2\ge C_1$.
\item[(III)] $\Lop[V]\ge2\mu\ve^{-2}q_n-\frac\mu4\ve^{-2}(c\Delta tK+q_n)-C_1\rho
\ge\frac{7\mu}4\ve^{-2}q_*-\frac{\mu cK}{4}\ve^{-2}\Delta t-C_1\rho
\ge\big(\tfrac{7\mu}{4}C_3-\tfrac{\mu cK}{4}-C_1\big)\rho\ge0$ if $C_3$ is large.
\item[(II)] $\Lop[V]\ge\sigma\theta-C_1\ve\rho-C_1\rho-C_1\ve^{-2}(\Delta t+q_n)
\ge\big(C_4\theta-C_1(2+1+C_3+C_2T)\big)\rho\ge0$ if $C_4$ is large.
\end{itemize}
Hence $V$ is a supersolution with $V^0_j\ge\Phi(x_j)=U^0_j$, and by
Lemma~\ref{lem:comp}
\[
U^n_j\le\Phi(x_j-ct_n+\sigma t_n)+q_n\le\Phi(x_j-ct_n)+KC_4T\rho+(C_3+C_2T)\ve^2\rho .
\]
This gives the first estimate. The second follows from
Lemma~\ref{lem:profile}(5),(1): $|\Phi_\ve(x-c_\ve t)-u(x,t)|\le K\ve+Ac_0|c_\ve-c_0|\,t\le C\ve$.
\end{proof}

\begin{remark}[What is lost]
The estimate is not sharp. Inside the reaction layer the truncation error is
pointwise of size $\Delta t/\ve^2$ (lagged reaction) and $h^2/\ve^3$
(second differences), and a pointwise barrier must absorb it with
$\sigma\Phi'$, so that the front is allowed to drift at speed $O(\rho)$.
The actual drift is the \emph{average} of the truncation error against the
adjoint eigenfunction of the linearized operator, which is a factor $\ve$
smaller: the experiments give speed errors $\approx-0.23\,\Delta t/\ve$ and
$\approx0.016\,(h/\ve)^2\cdot\ve^0$, i.e. $\Delta t/\ve$ and $(h/\ve)^2$.
\end{remark}

% ---------------------------------------------------------------------------
\section{Towards the sharp estimate}\label{sec:sharp}

The standard remedy is a corrector. Write the truncation error of (S) on
$\Phi$ as $\tau=\Delta t\,\tau_1+h^2\tau_2+\dots$ with $\tau_i$ explicit
functions of $z$, concentrated in the layer. Let $\mathcal A=\partial_z^2+c\partial_z-g_\ve'(\Phi)$
be the linearization at the wave. Since $\mathcal A\Phi'=0$, the equation
\[
\mathcal A\Psi_i=\tau_i+\lambda_i\Phi'
\]
is solvable (with bounded $\Psi_i$) iff
$\lambda_i=-\int\tau_i\psi^*/\int\Phi'\psi^*$, where $\psi^*=e^{cz}\Phi'$
spans the kernel of the adjoint. The barrier
$V=\Phi(z+\xi)+\Delta t\,\Psi_1(z+\xi)+h^2\Psi_2(z+\xi)\pm q$ with
$\xi'=\Delta t\lambda_1+h^2\lambda_2+$(higher order) has a residual one
order smaller, and the front drift is $\Delta t\lambda_1+h^2\lambda_2$. Scaling
gives $\lambda_1=O(1/\ve)$ and $\lambda_2=O(1/\ve^2)$, which are exactly the
observed $\Delta t/\ve$ and $(h/\ve)^2$. The technical points are: bounds on
$\Psi_i$ and its derivatives uniform in $\ve$ (in the inner variables this is a
regular problem), the behaviour of $\psi^*$ ahead of the front ($e^{cz}$
growth against the $e^{-cz}$ decay of $\Phi'$: $\psi^*$ is bounded), and the
nonsmoothness of $f$ at $1$ when only (F) holds.


\subsection*{Numerical check of the two mechanisms}
\texttt{numerics/experiments/truncation\_check.py} evaluates the truncation
error $\tau=\Lop[\Phi(x-ct)]$ of (S) on the exact wave ($f=(1-s)^+$, $A=1$),
averaged over the position of the front relative to the grid.
\begin{itemize}
\item \emph{Pointwise size} (what Theorem~\ref{thm:main} absorbs):
$\max|\tau|\approx a\,\Delta t/\ve^2+b\,h/\ve^2$ with $b\approx0.085$ ($w_0=0$)
and $b\approx0.47$ ($w_0=1$), for $\ve=0.01,0.005$ and $h=\ve/8,\ve/16$:
the $h/\ve^2$ scaling of the Lipschitz case.
\item \emph{Solvability (sharp) drift}
$\delta c=\int\tau\psi^*/\int\Phi'\psi^*$: for $w_0=0$ and $\Delta t=\kappa\ve^2$,
$\delta c/(\Delta t/\ve)=-0.2351,-0.2345,-0.2342$ for $\ve=0.02,0.01,0.005$
(independent of $\kappa$), in agreement with the speed error
$\approx-0.23\,\Delta t/\ve$ measured in the simulations of (S). The
spatial part of $\delta c$ is below $4\cdot10^{-5}$ for $h\le\ve/8$.
\end{itemize}
This supports the corrector programme of Section~\ref{sec:sharp}: the sharp
estimate should read $C_T(\ve+\Delta t/\ve+\text{(small spatial term)})$.

\medskip
\noindent\textbf{Plan.}
(1) Complete Lemma~\ref{lem:profile} (ODE phase-plane, uniform in $\sigma$).
(2) Numerical check of the barriers of Theorem~\ref{thm:main} (sign of $\Lop[V]$).
(3) Correctors and the sharp estimate.
(4) General data in 1D/radial with a classical solution: quantitative
version of Theorem~4.1 of the 2004 SIAM paper for the continuous problem, then the
discrete barriers as above.

\end{document}
